\chapter{모험심 노브}
% full version: chapters/08_nora.tex

점심 먹을 곳을 고른다. 늘 가던 식당에 갈까, 새 식당을 가 볼까? 늘 가던 곳만 가면 모퉁이 너머에 더 좋은 곳이 있다는 것을 영영 모른다. 늘 새 곳만 가면 맛없는 점심을 자주 먹는다. 알맞은 중간은 어디인가? 탐색할 것인가, 찾아낸 것을 활용할 것인가. 이 문제는 배우는 모든 존재 앞에 놓여 있다. 뇌에서는 노르아드레날린 방송국, \nt{NORA}가 이 문제를 맡는 것으로 보인다.

\section*{대비 노브}

노르아드레날린은 뉴런이 신호에 얼마나 날카롭게 응답하는지를 바꾼다. 작용이 약할 때 뉴런은 부드러운 밝기 조절기와 같다. 신호가 조금 커지면 응답도 조금 밝아진다. 작용이 강할 때는 스위치와 같다. 문턱값 아래에서는 침묵, 위에서는 최대 응답이다. 더 정확히 말하면, 실험은 노르아드레날린이 진짜 신호에 대한 응답보다 뉴런의 배경 잡담을 더 세게 억누른다는 것을 보여 준다. “대비 노브”는 이 효과를 다루기 편한 모델이다.

\pencilfig{8}{\pencilart{8}}{모험심 노브: 왼쪽으로 돌리면 기계가 새것을 시도하고, 오른쪽으로 돌리면 검증된 것을 단호하게 고른다.}

\section*{생각할 것인가 결정할 것인가}

3장에서 서로를 억제하는 두 무리를 떠올려 보자. 대비가 낮으면 둘 다 일하고 약한 쪽이 조금 조용할 뿐이다. 네트워크는 아직 “생각 중”이고 선택지를 저울질한다. 대비가 어떤 문턱값을 넘으면 네트워크는 갑자기 “결정함” 모드로 넘어간다. 한 무리가 이기고 다른 무리는 입을 다문다. 노브 하나가 기계를 숙고와 결정 사이에서 전환한다.

\section*{얼마나 탐색할까}

잡음을 더하자. 대비가 낮으면 잡음이 선택지 사이의 작은 차이를 쉽게 뒤집고, 기계는 최선이 아닌 선택지를 자주 시도한다. 탐색하는 것이다. 대비가 높으면 잡음은 아무것도 결정하지 못하고, 기계는 늘 최선이라고 여기는 것을 고른다. 활용하는 것이다.

우리 모델에서는 세상이 이따금 바뀐다. 최선의 선택지가 최악이 된다. 얻는 이득은 노브 위치에 따라 뒤집힌 U자 모양으로 변한다. 대비가 너무 적으면 기계는 눈먼 채 헤맨다. 너무 많으면 세상이 바뀐 것을 알아채지 못하고 맛이 변한 식당에 고집스레 계속 간다. 세상이 자주 바뀔수록 최적 설정은 낮아진다. 변덕스러운 세상에서는 탐색하는 편이 이득이다.

\pencilfig{108}{%
% points: models/nora_explore.py, reward as a share of the best possible, gain 0.5 ... 64 (doubling); the y axis starts at 0.70, hence the break mark
\draw[pen] (0,0) -- (6.3,0); \draw[pen] (0,0) -- (0,0.18); \draw[pen] (0,0.34) -- (0,2.4);
\draw[pcGray, line width=0.6pt] (-0.1,0.14) -- (0.1,0.22); \draw[pcGray, line width=0.6pt] (-0.1,0.3) -- (0.1,0.38);
\node[lbl, anchor=north] at (3.0,-0.05) {대비 노브};
\node[lbl, anchor=south, rotate=90] at (-0.1,1.3) {이득};
\draw[penlite=pcBlue] plot[smooth] coordinates {(0.00,0.55) (0.85,0.79) (1.70,1.19) (2.55,1.68) (3.40,2.00) (4.25,2.07) (5.10,2.01) (5.95,1.91)};
\draw[penlite=pcRed] plot[smooth] coordinates {(0.00,0.55) (0.85,0.69) (1.70,0.93) (2.55,1.24) (3.40,1.40) (4.25,1.28) (5.10,1.06) (5.95,0.89)};
\draw[dotpen=pcBlue, wash=pcBlue] (4.25,2.07) circle (0.07);
\draw[dotpen=pcRed, wash=pcRed] (3.40,1.40) circle (0.07);
\node[lbl, text=pcBlue, anchor=south] at (4.25,2.15) {세상이 드물게 바뀜};
\node[lbl, text=pcRed, anchor=north] at (3.40,1.30) {자주};
\node[lbl, anchor=north west, align=left] at (0.05,-0.35) {눈먼 채\\헤맨다}; \node[lbl, anchor=north east, align=right] at (6.3,-0.35) {변화를\\못 알아챈다};
}{노브 위치에 따른 모델의 이득. 곡선마다 꼭대기가 있다. 세상이 자주 바뀌면 꼭대기는 더 낮고 더 왼쪽에 있다.}

\section*{노브는 누가 돌리나}

변화의 자연스러운 징후는 놀람이다. 오차가 평소보다 커졌다는 것이다. 한 가설에 따르면 노르아드레날린이 알리는 것이 바로 이런 뜻밖의 변화이다. 여기서 미묘한 점이 중요하다. 일정한 배경 수준의 노르아드레날린은 오히려 탐색과 산만함과 관련되고, 중요한 사건에 대한 짧은 분출은 집중과 결정과 관련된다. 분출은 컴퓨터의 인터럽트처럼 작동하는지도 모른다. “하던 일을 멈춰라, 세상이 바뀌었다, 계획을 다시 짜라.”

이 시스템의 간접적인 징후는 바깥에서도 보인다. 무언가 뜻밖에 바뀌고 사람들이 더 빨리 배우기 시작할 때 동공이 커진다. 다만 동공은 노르아드레날린만 따르지 않으므로, 이것은 계측기가 아니라 단서이다.

그리고 솔직히 인정하자. 우리 모델에서 놀람에 따라 노브를 조정하는 방식은 가장 좋은 고정 설정보다 아주 조금만 더 나았다. 이런 조정이 정확히 어디서 제값을 하는지는 앞으로 밝혀야 한다.

\fullversion{8}
